\begin{scanpage}{0233}{215}
\begin{center}CONCLUSION\end{center}
In this book, the solution of a number of diffraction problems was obtained based on the approximate consideration of the field perturbation in the vicinity of a sharp bend of the surface or a sharp edge. Equations were derived for the scattering characteristics, or in certain cases (Chapter IV), for the radar reflection thicknesses with a specified irradiation direction. The expressions which have been found have a clear physical meaning. They satisfy the reciprocity principle, and they are convenient for making calculations.

The results obtained enable us to form a more complete concept of the applicability limits of the physical optics approach. It is usually customary to assume that this approach gives reliable results only if the body's dimensions are large in comparison with the wavelength. Such an opinion is based on the following argument. The physical optics approach considers only the radiation from the uniform part of the current, and does not include in the calculations the nonuniform part of the current which is concentrated in the vicinity of the bends and the sharp edges. Therefore, when the body's dimensions are considerably larger than the wavelength, the nonuniform part of the current occupies a relatively small part of the body's surface. Therefore, one would think that its influence would be small.

But in actuality it turns out that the reliability of physical optics results depends substantially, not only on the body's dimensions, but also on the body's shape and the irradiation and observation directions. For example, with the glancing incidence of a wave on the flat face of a body, the edge zone occupied by the nonuniform part of the current is considerably broadened and the effect of this current becomes substantial. Therefore, physical optics gives qualitatively incorrect results for the field scattered by flat plates with glancing irradiation independent of the ratio between their dimensions and the wavelength. The effect of the nonuniform part of
\end{scanpage}
